\chapter{The neural network and its training, step by step}\label{ch:network}
\paperloc{Sections 2B--2C, Fig.~4, and the reported network architecture.}

\section{A neuron is a weighted sum followed by a function}
Given inputs $a_1,\ldots,a_d$, a neuron forms
\[
q=\sum_{j=1}^dw_ja_j+b,\qquad a_{\rm out}=g(q).
\]
The weights $w_j$ determine which inputs contribute and with what sign;
the bias $b$ shifts the response. The activation $g$ introduces nonlinearity.
A layer of neurons applies this operation in parallel:
\begin{equation}
\vct q^{(\ell)}=A^{(\ell)}\vct a^{(\ell-1)}+\vct b^{(\ell)},
\qquad
\vct a^{(\ell)}=g_\ell(\vct q^{(\ell)}).
\label{eq:network}
\end{equation}
The function acts component by component. Repeating layers produces a
composition of nonlinear functions. With only linear activations,
all layers collapse into a single affine map, no matter how many are used.

For example, $g(q)=\tanh q$ is smooth, ranges from $-1$ to $+1$,
and has derivative $1-\tanh^2q$. A rectified linear unit is
$\max(0,q)$ and is piecewise linear. The paper does not report its
activation functions, so neither choice can be attributed to the authors.
Likewise, an ordinary unbounded linear regression output is a possible
choice, but the actual output activation is undocumented.

\section{What the network is asked to learn}
Here the input is the 84-component imaging vector, not a pupil image.
The output is a vector of Zernike coefficients, not CD or PVB.
The desired relationship is
\[
G_\Theta:\vct y\longmapsto\vct c,
\]
where $\Theta$ collects all trainable weights and biases.
During training, the known simulator pair
$(\vct y^{(r)},\vct c^{(r)})$ supplies the input and target.
During use, a new desired imaging vector supplies the input.

Calling this a backpropagation network describes how derivatives of the
training loss are computed. It does not mean light is physically
propagated backward through a Maxwell solver.
Back prediction refers to the direction of the learned inverse mapping;
backpropagation refers to the chain rule used to train it.

\section{Interpreting the reported layer widths}
The paper gives
\[
84,\ 84,\ 64,\ 64,\ 50,\ 36,\ 32,\ 25,\ 25.
\]
Read literally, the network has 84 inputs and 25 outputs. Its seven
hidden layers have widths $84,64,64,50,36,32,25$.
Figure 4 is a simplified schematic and does not depict these exact widths
or the full hidden-layer count.

A fully connected layer taking $d_{\rm in}$ inputs to $d_{\rm out}$
outputs has $d_{\rm in}d_{\rm out}$ weights and $d_{\rm out}$ biases.
The literal architecture therefore has
\[
\sum_{\ell=1}^8(d_{\ell-1}+1)d_\ell
=7140+5440+4160+3250+1836+1184+825+650
=24\,485
\]
trainable parameters. This calculation assumes the conventional dense
architecture; the publication does not provide implementation code to
confirm all details.

The parameter count exceeds the number of simulated aberration cases.
That does not mathematically prove overfitting: each case has multiple
outputs, and training structure matters. It does make independent testing,
preprocessing disclosure, and regularization choices important.
The 25-output ambiguity discussed in Chapter \ref{ch:data} remains.

\section{Why feature normalization matters}
CD may be of order tens of nanometres while $\mathrm{PS}_y$ can be a small
fraction of a nanometre. An optimizer acting on raw coordinates may assign
very different numerical scales to these quantities.
Featurewise min--max normalization to $[-1,1]$ is
\begin{equation}
\widetilde x=2\frac{x-x_{\min}}{x_{\max}-x_{\min}}-1.
\label{eq:normalization}
\end{equation}
Solving for the physical value,
\begin{equation}
x=x_{\min}+\frac{\widetilde x+1}{2}(x_{\max}-x_{\min}).
\label{eq:unnormalization}
\end{equation}
The paper names MATLAB's \texttt{mapminmax}, which performs rowwise
range mapping and supports applying stored settings to new data
\cite{mapminmax}.

The minimum and maximum must be fitted on training data and saved.
Use those same settings for validation, testing, and deployment.
Recomputing them independently for a new batch changes the meaning
of every network coordinate. Including held-out data when fitting them
lets test information enter preprocessing.

If $x_{\max}=x_{\min}$, the feature has no training variation and the
formula divides by zero. It needs explicit handling.
If a new value lies outside the training range, its normalized value
lies outside $[-1,1]$; min--max mapping does not by itself clip it.
Similarly, undoing output normalization does not impose a coefficient
bound unless the output is explicitly bounded or clipped.

\section{The training loss}
Let $N$ be the number of training samples and $K$ the number of output
coordinates actually trained. One common mean-square loss is
\begin{equation}
L(\Theta)=\frac{1}{NK}\sum_{r=1}^N\sum_{j=1}^K
\left[\widehat c_j^{(r)}-c_j^{(r)}\right]^2.
\label{eq:trainingMSE}
\end{equation}
Often the coefficients here are normalized coordinates.
The paper says training uses MSE and a goal of $10^{-8}$, but does not
supply enough preprocessing detail to turn that number into a guaranteed
physical coefficient error.

If, purely as an example, one output spans
$[-20,+20]\,\mathrm{m}\lambda$ and is mapped to $[-1,1]$, then
a normalized RMSE $10^{-4}$ corresponds to
$0.002\,\mathrm{m}\lambda$ for that coordinate.
The square root of an average training MSE does not bound each output,
does not describe held-out error, and does not bound image error after
the nonlinear forward map.

\section{Backpropagation is repeated application of the chain rule}
For clarity use a one-sample loss
$L=\tfrac12\norm{\widehat{\vct c}-\vct c}^2$.
For a linear output layer, the derivative with respect to its preactivation is
\[
\vct\delta^{(L)}=\widehat{\vct c}-\vct c.
\]
For a nonlinear output, multiply componentwise by $g_L'$.
At each earlier layer,
\begin{equation}
\vct\delta^{(\ell)}
=\left[(A^{(\ell+1)})^T\vct\delta^{(\ell+1)}\right]
\odot g_\ell'(\vct q^{(\ell)}).
\label{eq:backprop}
\end{equation}
The symbol $\odot$ means componentwise multiplication.
To understand the transpose, note that one hidden neuron influences
every neuron in the next layer; its derivative sums all those downstream
paths, weighted by the connecting weights.

The parameter derivatives are
\begin{equation}
\frac{\partial L}{\partial A^{(\ell)}}=
\vct\delta^{(\ell)}(\vct a^{(\ell-1)})^T,\qquad
\frac{\partial L}{\partial\vct b^{(\ell)}}=\vct\delta^{(\ell)}.
\label{eq:networkgrad}
\end{equation}
Average or sum over samples as required by the loss normalization.
Backpropagation computes these derivatives efficiently by reusing
intermediate derivatives, rather than perturbing every weight separately.

\worked{One hidden neuron}
Take $a=\tanh(wx+b)$, $\hat c=va+d$, and
$L=\tfrac12(\hat c-c)^2$.
With $e=\hat c-c$,
\[
\frac{\partial L}{\partial v}=ea,\quad
\frac{\partial L}{\partial d}=e,\quad
\frac{\partial L}{\partial w}=ev(1-a^2)x,\quad
\frac{\partial L}{\partial b}=ev(1-a^2).
\]
Every factor has a role: residual, downstream weight, activation slope,
and upstream input. A deep network repeats the same multiplication and
summation pattern across more paths.

\section{Deriving the Levenberg--Marquardt update}
Flatten all training residuals into a vector $\vct e(\Theta)$ and let
$J_e=\partial\vct e/\partial\Theta$.
This \emph{training Jacobian} differentiates residuals with respect to
network parameters. It is distinct from the imaging Jacobian
$\partial\vct y/\partial\vct c$.

For a small parameter change $\Delta\Theta$,
$\vct e(\Theta+\Delta\Theta)\simeq\vct e+J_e\Delta\Theta$.
Minimize the locally linearized residual with a damping penalty:
\[
\frac12\norm{\vct e+J_e\Delta\Theta}^2
+\frac{\mu}{2}\norm{\Delta\Theta}^2.
\]
Differentiate with respect to $\Delta\Theta$, set the derivative to zero,
and obtain
\begin{equation}
(J_e^TJ_e+\mu I)\Delta\Theta=-J_e^T\vct e.
\label{eq:LM}
\end{equation}
Solve this linear system; do not explicitly form a matrix inverse
in a numerical implementation.

For small $\mu$, the step approaches Gauss--Newton.
For very large $\mu$,
$\Delta\Theta\simeq-\mu^{-1}J_e^T\vct e$,
which is a small gradient-descent step.
The optimizer adjusts damping according to whether a proposed step
improves the loss. This is the mathematical reason LM can combine
rapid local convergence with caution away from a good fit.

For a squared-residual loss, the exact Hessian contains
$J_e^TJ_e+\sum_i e_i\nabla^2e_i$.
Gauss--Newton drops the second term, which is often small near a
good fit. LM adds the stabilizing diagonal term.
It does not guarantee finding a global optimum of a deep network.

The paper identifies MATLAB's LM training function. Its official
documentation distinguishes the MSE performance goal from
validation-based stopping and documents the damping update
\cite{trainlm}. These software facts confirm the interpretation of the
algorithm; the derivation above follows directly from the local
least-squares problem.

\section{Training, validation, and test have different jobs}
Training data determine the weights. Validation data guide choices such
as stopping time or architecture. Test data estimate performance after
those choices are fixed. Repeatedly changing the architecture after
looking at the same test set makes that set part of model selection.

The paper says training stops below MSE $10^{-8}$ or after six rounds
without an obvious decrease. Standard MATLAB \texttt{trainlm} uses a
six-failure default for \emph{validation} performance.
The paper's wording does not settle whether its implemented condition
was precisely this default or a different training criterion.
Record that ambiguity rather than equating training stagnation with
validation failure.

\section{What the learning theory does and does not promise}
The Introduction cites approximation results for neural networks.
Such results motivate using nonlinear networks to approximate continuous
functions under suitable conditions. They do not establish that this
particular architecture, finite dataset, and optimizer will recover a
stable inverse of an ambiguous optical map.

In particular, they do not guarantee extrapolation beyond the training
region, compliance with coefficient bounds, a manufacturable optical
system, or reliable behavior when the source or mask changes.
The paper's useful evidence is its forward-simulation validation,
not the general existence of approximation theorems.

\chapter{Errors, percentages, and what validation means}\label{ch:errors}
\paperloc{Equations (2)--(3), Figs.~5--10, and the statements about small PS values.}

\section{Root mean square error}
For predictions $\hat y_i$ and reference values $y_i$, define
$e_i=\hat y_i-y_i$. Then
\begin{equation}
\mathrm{RMSE}=\sqrt{\frac1N\sum_{i=1}^Ne_i^2}.
\label{eq:RMSE}
\end{equation}
This is the paper's Eq.~(2). Squaring prevents positive and negative
errors from cancelling. Taking the square root restores the units.
RMSE of a CD in nm is in nm. RMSE of a coefficient in waves is in waves.
The figure axis must be checked before comparing magnitudes.

If $\bar e$ is the error mean and
$s_e^2=N^{-1}\sum_i(e_i-\bar e)^2$, direct expansion gives
\[
\mathrm{RMSE}^2=\bar e^2+s_e^2.
\]
RMSE combines systematic bias and scatter.
It is not the maximum error. For example, 99 exact predictions and one
1 nm error have RMSE 0.1 nm over 100 cases.

\section{Mean absolute percentage error}
The paper's Eq.~(3) is
\begin{equation}
\mathrm{MAPE}
=\frac{100\%}{N}\sum_{i=1}^N
\left|\frac{\hat y_i-y_i}{y_i}\right|.
\label{eq:MAPE}
\end{equation}
The absolute value makes each term nonnegative, including for negative
reference values. The denominator makes the statistic dimensionless.
It also gives large weight to small reference magnitudes and is undefined
when $y_i=0$.

For a 14 nm CD, an error of 0.014 nm is 0.1 percent.
For a 0.01 nm pattern shift, the same 0.014 nm error is 140 percent.
The second percentage looks alarming, but the physical displacement
error is identical. Conversely, a 1 percent error on a large quantity
can exceed a stringent absolute tolerance.

\figteach{percentage_error}{1.0}{A fixed absolute error of 0.01 nm
produces an arbitrarily large percentage as the reference shift approaches
zero. The right panel shows that the physical error does not diverge.}

\section{Why near-zero shifts are especially awkward}
Zero is often the desired shift. A metric that becomes undefined exactly
at the target is a poor sole measure of control quality.
The paper separates values around 0.14 nm in Fig.~8 and around 0.1 nm
in Fig.~10. The 0.14 nm scale is $1\%$ of the nominal 14 nm width;
it is not the paper's final $\pm0.4\nm$ acceptance boundary.

The text describes values above and below these thresholds but does not
fully document how signed shifts, exact zeros, or removed values were
handled. A careful replication would define whether the split uses
$|y_i|$, explicitly state any denominator floor, and report the sample
count in each group.

A worsening MAPE for smaller $|y|$ can occur with unchanged absolute
error. Therefore the graph alone does not demonstrate that a neural
network intrinsically cannot predict small numbers.
Noise, feature scaling, cancellation between coefficients, weak
observability, target density, and preprocessing can all contribute.
The paper's explanation should be treated as a suggested account,
not a mathematical inevitability of neural networks.

\section{Better companion measures for a placement budget}
Keep RMSE and add interpretable quantities:
\begin{align}
\mathrm{MAE}&=\frac1N\sum_i|e_i|,\\
e_{\max}&=\max_i|e_i|,\\
\mathrm{pass\ fraction}&=
\frac{\#\{i:|\mathrm{PS}_{y,i}^{\rm verified}|\leq b\}}{N}.
\end{align}
For a specified tolerance $b>0$, a useful normalized error is
$|e_i|/b$. Its denominator remains meaningful at zero target shift.
Quantiles show tail behavior; signed mean error detects systematic bias.
None of these removes the need to state which patterns and operating
conditions are included.

\worked{A candidate near the boundary}
Suppose desired shift is 0.395 nm and the prediction-to-simulation
error is $+0.020\nm$. The verified shift is 0.415 nm and fails a
0.4 nm limit. A small absolute discrepancy is consequential here.
The same discrepancy at a target of zero would still meet the limit.
The relevant question is error relative to remaining margin, not just
whether a number sounds small.

\section{Coefficient error versus image error}
For a held-out original pair $(\vct c,\vct y=F(\vct c))$,
the network outputs $\hat{\vct c}=G(\vct y)$.
Two separate residuals are
\[
\vct r_c=\hat{\vct c}-\vct c,\qquad
\vct r_y=F(\hat{\vct c})-\vct y.
\]
Figure 5 concerns coefficient differences. Figures 6 and 7 concern the
imaging differences after forward simulation.
A sizeable $\vct r_c$ can coexist with a small $\vct r_y$ if the two
wavefronts are nearly equivalent for the selected indicators.
However, coefficient disagreement by itself does not prove a valid
alternative: the forward residual must actually be checked.

Agreement in a few scalar metrics also does not establish identical
full contours, NILS, stochastic defect behavior, or performance on
unseen pattern classes. The paper validates the indicators it measures.

\section{What an independent simulation test establishes}
The 118 additional target vectors provide a useful test beyond the model's
training/validation set, as described by the authors.
Their proposed coefficients are sent through the same type of forward
simulator under the same lithographic conditions.
This tests consistency with the simulated model and conditions.
It is not an experimental measurement on a fabricated projector or wafer.

Using the same simulator for training labels and final checking can
validate the inverse approximation while leaving a common forward-model
bias invisible. Actual engineering deployment also needs calibrated
mask, resist, optical, and metrology models.
These limitations define the scope of the evidence rather than negate
the usefulness of a simulation-based study.
